\cotefile{78}
\scanpage{32}
\deleted{et enfin la condition d'unimodularité}

\deleted{$4^\circ)$ $u_E$ unimodulaire}

% vertical rule beside the definition
On dit que la \deleted{\unsure{repr.}}\inserted{réalisation} régulière est \uline{unimodulaire} si elle satisfait
\[
4^\circ)\quad u_E \text{ est unimodulaire}
\]

\uline{Scholie} En \unsure{étudiant|étendant} ces objets, on étudie simultanément tous les $n$-polygônes réguliers \inserted{épinglés} pour $n\in\unsure{\mathbb N'|\mathbb N|\mathbb N^{*}},\ n\geq 3$ -- ils correspondent aux cas où $u_E$ est d'ordre fini, et $n$ est l'ordre de $u_E$ -- \unsure{en|et} même temps qu'un nouveau type d'objets : les $\infty$-polygônes réguliers épinglés.

\uline{Remarque} Les considérations générales concernant les ``polygônes réguliers'' s'étendent aussi aux $\infty$-polygônes réguliers sans chgt. ! Notamment les prop. 1, et corollaire, la prop 2 (qui \deleted{\unsure{est}}\inserted{s'étend en l'énoncé évident} concernant la \deleted{du} question de fidélité $D_\infty\hookrightarrow\mathrm{Aff}(E)$ exprimant \unsure{qu'on|que} \deleted{\unsure{conditions que|conditions}} \unsure{l'on.|bien.|finies} bien un $\infty$-poly. rég-
% slanted
\margin{! Mais il faut commencer par donner la définition des polyèdres réguliers en général (finis ou pas), en \deleted{\unsure{répét.}}\unsure{reprenant} la définition d'abord donnée dans le cas fini \ldots}
% archive sheet number: 31
